\documentclass{article}


% pacotes
\usepackage[brazil]{babel}
\usepackage[letterpaper,top=2cm,bottom=2cm,left=3cm,right=3cm,marginparwidth=1.75cm]{geometry}
\usepackage{fancyhdr}

\pagestyle{fancy}       % Ativa o estilo fancyhdr
\fancyhf{}              % Limpa todos os cabeçalhos e rodapés padrão
\renewcommand{\headrulewidth}{0pt}

% rodapé
\fancyfoot[L]{Relatório}
\fancyfoot[R]{\thepage}



\begin{document}
%capa
\begin{titlepage}
    \centering
    \vspace*{2cm}
    {\large
    Beatriz São Leandro Cosimatti\\
    17024130\\
    João Victor Bozola Bosi\\
    5979703\\
    João Vitor Vasconcelos Jacob\\
    nUsp do jacob\\
    Laís Antunes Cayres Quintão\\
    16880392\\ 
    }
    \vspace{2cm}
    \vspace{3cm}
    
    {\Large {Trabalho de Estrutura de Dados II:\\}}
    {\Large \textbf{Análise e Implementação de Algoritmos\\}}
    \vspace{3cm}
    {\large
    Disciplina: Estrutura de Dados II\\
    SCC0606\\
    Professor: 
    }
    \vfill
    {\large
    São Carlos\\
    {14 jun. 2026}
    }
\end{titlepage}


\section{Introdução}
intro\\
Acho que coloca o link do vídeo aqui

\newpage

\section{Algoritmos Base}

    \subsection{Bubble Sort}
    explicação do Bubble


    \subsection{Insertion Sort}
    explicação do Insertion


    \subsection{Quick Sort}
        O algoritmo de ordenação Quick Sort foi desenvolvido e publicado pela primeira vez em 1961, pelo cientista da computação britânico Charles Antony Richard Hoare (1934 - 2026). A técnica de ordenação foi concebida enquanto o cientista estudava na Universidade de Moscou e trabalhava no \textit{National Physical Laboratory (NPL)}, do Reino Unido. Na ocasião, Hoare estava participando de um projeto que tinha como objetivo a tradução de textos em russo para inglês e, para isso, como os dicionários erar armazenados em fita magnética, necessitava ordenar das palavras dentro de uma sentença em ordem alfabética. Nesse contexto, surgiu a necessidade de desenvolver um algoritmo de ordenação que fosse assintoticamente menor que $O(n^2)$.

    \subsection{Heap Sort}
    explicação do Heap


    \subsection{Radix Sort}
    explicação do radix

\newpage

\section{Arquitetura do Sistema}
Descrição:
• das métricas coletadas;
• das heurísticas;
• das decisões implementadas -> parte da bia

\newpage

\section{Metodologia Experimental}
Descrever:
• ambiente utilizado;
• datasets;
• tamanhos testados;
• critérios de avaliação

\newpage

\section{Resultados}
Apresentar tabelas, gráficos, métricas coletadas, com discussão crítica. Explicar:
• decisões corretas;
• decisões ruins;
• comportamentos inesperados;
• limitações observadas;
• impacto das características da entrada.
• impactos das entradas adversariais

\newpage

\section{Reflexão sobre Uso de IA}
O grupo deverá informar:
• quais ferramentas de IA foram utilizadas;
• em quais etapas auxiliaram;
• erros encontrados nas respostas geradas;
• correções realizadas pelo grupo.

\newpage

\section{Conclusão}
Resumo dos principais achados.

\newpage

\section{Referências Bibliográficas}

\end{document}
